\part{Architecture}\label{ds:part:architecture}

\chapter{System overview}\label{ds:ch:system-overview}

This chapter gives the top-level view of the system. It lists the major blocks and shows how they connect.

\chapter{CPU phases 0 to 3}\label{ds:ch:cpu-phases-0-to-3}

This chapter states the requirements for the CPU in phases 0 to 3. It covers the core and the cache system.

\chapter{GPU-lite SIMT engine}\label{ds:ch:gpu-lite-simt-engine}

This chapter states the architecture of the GPU-lite engine. It defines the instruction set and the execution model.

\chapter{SoC integration}\label{ds:ch:soc-integration}

This chapter states how the blocks join in the SoC. It defines the bus fabric and the integration rules.

\chapter{Phase 6 IP expansion}\label{ds:ch:phase-6-ip-expansion}

This chapter states the blocks that Phase 6 adds. It lists the peripherals, the crypto accelerator and the NPU.

\chapter{Mixed-signal PLL}\label{ds:ch:mixed-signal-pll}

This chapter states the requirements for the PLL. It covers the analog design and the model that the SoC uses.

\chapter{Memory map}\label{ds:ch:memory-map}

This chapter gives the address space. It lists the register map of each block.

\chapter{Interface definitions}\label{ds:ch:interface-definitions}

This chapter defines the signals of each module interface.

\chapter{Reference models}\label{ds:ch:reference-models}

This chapter describes the Python reference models. It defines the model API.

\chapter{Locked architecture decisions}\label{ds:ch:locked-architecture-decisions}

This chapter lists the architecture decisions that stay fixed. It gives the reason for each decision.

% Moved verbatim from CLAUDE.md by tools/docs/md2tex.py; edit this file, not the original.
\chapter{Key Design Decisions (From Specifications)}\label{ds:dec:key-design-decisions}
\begin{relocated}{CLAUDE.md}

\textbf{Phase 1 CPU}:

\begin{itemize}
\tightlist
\item
  Single-cycle execution with AXI stalls for memory operations
\item
  Unified AXI4-Lite bus for instruction fetch and data access
\item
  x0 register hardwired to zero (per RV32I spec)
\item
  EBREAK instruction triggers CPU halt; debug writes only allowed when
  halted
\item
  Only illegal instruction traps supported (no interrupts in Phase 1)
\item
  Naturally aligned memory accesses only (misaligned = trap)
\end{itemize}

\textbf{Phase 4 GPU}:

\begin{itemize}
\tightlist
\item
  SIMT execution model (8 lanes per warp); single compute unit
\item
  Round-robin warp scheduling; one-level divergence handling
\item
  No cache coherence with CPU; memory coalescing when possible
\end{itemize}

\textbf{Sky130 frequency ceiling}: \textasciitilde100--120 MHz maximum
for this CPU pipeline topology (see Technology Node Strategy). PDK
limitation --- not an RTL issue. For higher frequencies, target
FreePDK45/\allowbreak{}NanGate45 (150--250 MHz) or ASAP7 (\textbf{1418
MHz demonstrated} at Run 43, 2026-05-20).
\end{relocated}


% Moved verbatim from CLAUDE.md by tools/docs/md2tex.py; edit this file, not the original.
\chapter{Phase 5 SoC locked architecture decisions}\label{ds:dec:locked-phase5}
\begin{relocated}{CLAUDE.md}

\textbf{Phase 5 SoC --- Locked Architecture Decisions} (confirmed
2026-03-28):

\begin{enumerate}
\def\labelenumi{\arabic{enumi}.}
\tightlist
\item
  \textbf{Cache line size}: 16 bytes locked for Phase 3. Revisit 32-byte
  lines only if/when an L2 cache is added.
\item
  \textbf{Cache associativity}: Direct-mapped locked. \textbf{M10
  RESOLVED (2026-06-21): stays direct-mapped.} A 2-way D\$ was
  prototyped + benchmarked (\code{experiment/dcache-2way}): it
  eliminates synthetic K=2 conflict misses (153/\allowbreak{}kI→0) but
  representative firmware (CoreMark, 232 B working set) fits L1 with 18×
  margin and shows no conflict pressure --- so 2-way is not adopted
  (adds \textasciitilde100--200 ps hit-path mux for no real benefit).
  See \src{docs/M10_L2_DECISION_ANALYSIS.md}.
\item
  \textbf{AXI4 burst upgrade}: Phase 3 refill FSMs use 4 sequential
  AXI4-Lite transactions. These are upgraded to AXI4 burst mode during
  Phase 5 SoC integration --- do not change Phase 3 FSMs before then.
\item
  \textbf{L2 cache}: Not in Phase 5 scope. \textbf{M10 RESOLVED
  (2026-06-21): NO-GO.} Benchmarking (\src{sw/bench/}
  sweep/\allowbreak{}conflict/\allowbreak{}CoreMark on \code{soc_top})
  found the L1 capacity wall (\textgreater4 KB) is synthetic-only;
  representative firmware fits the 4 KB L1; the GPU carries bulk data.
  No real L1 bottleneck → no L2. Revisit only if a future on-core
  workload sustains \textgreater4 KB working sets. See
  \src{docs/M10_L2_DECISION_ANALYSIS.md}.
\item
  \textbf{DRAM controller}: Use behavioral AXI4-slave SRAM model for
  Phase 5. Real DRAM controller (with refresh) is a Phase 6+ stretch
  goal only if tape-out is targeted.
\end{enumerate}
\end{relocated}


% Moved verbatim from CLAUDE.md by tools/docs/md2tex.py; edit this file, not the original.
\chapter{Phase 4: GPU-Lite SIMT Engine (frozen)}\label{ds:ch:phase4-5-architecture}
\begin{relocated}{CLAUDE.md}

\begin{itemize}
\tightlist
\item
  \textbf{ISA}: Custom 32-bit vector encoding --- VADD, VSUB, VMUL,
  VAND, VOR, VSLL, VLD, VST, VBEQ, VBNE, VJMP, VRET, VSYNC
\item
  \textbf{Execution}: Single compute unit; 8 SIMD lanes; 8 warps max; 64
  total threads
\item
  \textbf{Register file}: 32 registers × 8 threads per warp (1 KB)
\item
  \textbf{Divergence}: Single-level SIMT stack with per-lane active
  mask; round-robin scheduler
\item
  \textbf{Shared memory}: 16 KB scratchpad, 32 banks
\item
  \textbf{Memory}: AXI4 master (burst) + memory coalescer (8-lane →
  fewer AXI4 bursts)
\item
  \textbf{Coherency}: Software-managed --- CPU flushes D-cache before
  launch, invalidates after completion. No HW coherence.
\item
  \textbf{Spec}: \src{docs/design/PHASE4_GPU_ARCHITECTURE_SPEC.md}
\end{itemize}

\section{\texorpdfstring{Phase 5: SoC Integration (\ok{} COMPLETE
2026-06-24)}{Phase 5: SoC Integration ( COMPLETE 2026-06-24)}}\label{ds:ch:phase4-5-architecture:phase-5-soc-integration-white-check-mark-complete-2026-06-24}

\begin{itemize}
\tightlist
\item
  \textbf{Interconnect}: AXI4 crossbar (data fabric) + AXI-Lite bus
  (control/\allowbreak{}config)

  \begin{itemize}
  \tightlist
  \item
    AXI4 masters: CPU D-cache (upgraded to burst), GPU, DMA engine
  \item
    AXI4 slave: behavioral SRAM controller (no DRAM refresh)
  \item
    AXI-Lite slaves: GPU/DMA config registers, UART, SPI, timer, IRQ
    controller
  \end{itemize}
\item
  \textbf{Cache upgrade}: Phase 3 refill FSMs upgraded from 4 sequential
  AXI4-Lite beats to AXI4 burst transactions
\item
  \textbf{DMA engine}: Descriptor queue, AXI4 master, interrupt output
  to CPU
\item
  \textbf{Memory model}: Behavioral AXI4-slave SRAM (no DRAM refresh
  logic)
\item
  \textbf{Power domains}: PD\_CPU, PD\_GPU, PD\_SRAM, PD\_PERIPH (clock
  + power gating per \code{UPF_POWER_SPEC.md})
\item
  \textbf{Performance counters}: CSR-mapped --- cycle count,
  instructions retired, branch mispredictions, I\(/D\) miss counts,
  active warps, warp stall cycles, divergence events
\item
  \textbf{Roadmap}: \src{docs/PHASE5_SOC_INTEGRATION_PLAN.md} (12
  milestones M1--M12, golden spec)
\end{itemize}
\end{relocated}


% Moved verbatim from CLAUDE.md by tools/docs/md2tex.py; edit this file, not the original.
\chapter{Struct-member select rule}\label{ds:inv:struct-member-select}
\begin{relocated}{CLAUDE.md}

\textbf{Struct-member select rule (\prref{275}, bead \bead{ainf}; enforced
by a blocking \code{qa-checks} gate)}: no bit, part or element select
may be taken directly on a struct member anywhere in \src{rtl/}, read or
write. Copy the member to a named wire first. The checker is
\src{tools/verif/check_struct_member_select.py}, and its waiver file
\src{tools/verif/struct_member_select_baseline.txt} must stay empty.
Synlig elaborates this idiom silently and wrongly
(\src{docs/SYNLIG_STRUCT_SELECT_AINF.md}); the \code{ma7}/\code{dud4}
entries above trace to the same defect class.
\end{relocated}


% Moved verbatim from CLAUDE.md by tools/docs/md2tex.py; edit this file, not the original.
\chapter{Strict elaboration rule}\label{ds:inv:strict-elaboration}
\begin{relocated}{CLAUDE.md}

\textbf{Strict elaboration rule (bead \bead{lvdk}; blocking
\code{qa-checks} gate)}: \src{tools/verif/check_elab_strict.py} runs
pyslang over \code{soc_top}, \code{rv32i_cpu_top} and \code{gpu_top} on
6 arms (3 default, FPGA, 2 macro), because Verilator and sv2v accept
use-before-declare while VCS and Xcelium reject it. Add a new define arm
to its \code{ARMS} table. Only macro arms may use
\code{--ignore-unknown-modules}, which also hides genuinely missing
modules.
\end{relocated}
